% !TeX program = xelatex
\documentclass[12pt,a4paper]{article}
\usepackage{fontspec}
\setmainfont{Liberation Serif}   % 英文主字体
\usepackage{xeCJK}
\setCJKmainfont{Microsoft YaHei} % 中文主字体（可替换为 SimSun, Noto Sans CJK SC）
\setCJKmonofont{Microsoft YaHei}

\usepackage[english]{babel}      % 保留英文环境，不影响中文
\usepackage{geometry}
\geometry{margin=1in}
\usepackage{amsmath,amssymb}
\usepackage{booktabs}
\usepackage{longtable}
\usepackage{array}
\usepackage{xcolor, colortbl}      % 用于 \rowcolor
\usepackage{hyperref}

\hypersetup{
	colorlinks=true,
	linkcolor=blue,
	citecolor=blue,
	urlcolor=blue,
}

\title{协调物质周期拉格朗日量（PLCM）的补充校准数据}
\author{阿列克谢·V·雅库舍夫}
\date{2026年3月}

\begin{document}
	\maketitle
	
	\tableofcontents
	\newpage
	
	\section{引言}
	本文档提供了协调物质周期拉格朗日量（PLCM）框架中使用的完整校准数据集。所有值均源自主要 PLCM 附录中描述的公式和方法，并以机器可读格式存储在存储库 \url{https://github.com/Alexey-Yakushev-YUCT/YPSDC/} 中。下面的表格给出了简要视图；完整的 $118\times118$ 矩阵和全面的二元数据在存储库中作为 CSV 文件提供。
	
	\section{基线耦合系数 $\kappa_{ij}^{\text{baseline}}$}
	基线耦合系数使用以下公式计算：
	\[
	\kappa_{ij}^{\text{baseline}} = \frac{2\,\Delta H_{\text{mix}}^{ij}}{\Delta H_{\text{mix,ref}}} \cdot
	\frac{1}{1 + |r_i - r_j|/r_{\text{avg}}} \cdot
	\frac{1}{1 + |\chi_i - \chi_j|},
	\]
	其中 $\Delta H_{\text{mix,ref}} = -7.5$ kJ/mol（Cu–Sn 参考）。对于缺乏实验混合焓的系统，使用 Miedema 模型 \cite{Miedema1980}。
	
	完整的 $118\times118$ 矩阵是对称的，在存储库中作为 \texttt{kappa\_baseline.csv} 提供。前 20 个元素的代表性片段在主要 PLCM 文档中给出。
	
	\section{特定于性质的校准因子 $\beta_{ij}^{(P)}$}
	对于每个二元系统 $i$–$j$ 和性质 $P$，耦合系数按因子 $\beta_{ij}^{(P)}$ 缩放：
	\[
	\kappa_{ij}^{(P)} = \beta_{ij}^{(P)} \cdot \kappa_{ij}^{\text{baseline}}.
	\]
	
	表 \ref{tab:beta_group} 列出了相似元素组的平均 $\beta$ 值，而表 \ref{tab:beta_special} 给出了选定的高影响力系统（包括新的 Mg–Zn 数据）的值。
	
	\begin{longtable}{p{2.5cm} p{2.5cm} p{1.2cm} p{2.5cm} p{2.5cm} p{1.2cm}}
		\caption{按元素组的校准因子 $\beta_{ij}^{(P)}$（所有性质）}\label{tab:beta_group}\\
		\toprule
		\textbf{组} & \textbf{系统} & $\beta$ & \textbf{组} & \textbf{系统} & $\beta$ \\
		\midrule
		\endfirsthead
		\toprule
		\textbf{组} & \textbf{系统} & $\beta$ & \textbf{组} & \textbf{系统} & $\beta$ \\
		\midrule
		\endhead
		\multicolumn{6}{c}{\textbf{碱金属（Li, Na, K, Rb, Cs, Fr）}} \\
		Li–Na & 所有 & 0.50 & K–Rb & 所有 & 0.55 \\
		Li–K  & 所有 & 0.48 & Rb–Cs & 所有 & 0.52 \\
		\multicolumn{6}{c}{\textbf{碱土金属（Be, Mg, Ca, Sr, Ba, Ra）}} \\
		Be–Mg & 所有 & 0.60 & Mg–Ca & 所有 & 0.55 \\
		Ca–Sr & 所有 & 0.58 & Sr–Ba & 所有 & 0.56 \\
		\multicolumn{6}{c}{\textbf{过渡金属（第 3–12 族）}} \\
		Sc–Ti & 所有 & 0.75 & Ti–V & 所有 & 0.80 \\
		V–Cr  & 所有 & 0.85 & Cr–Mn & 所有 & 0.82 \\
		Mn–Fe & 所有 & 0.88 & Fe–Co & 所有 & 0.90 \\
		Co–Ni & 所有 & 0.92 & Ni–Cu & 所有 & 0.88 \\
		Cu–Zn & 所有 & 0.80 & Zn–Ga & 所有 & 0.70 \\
		Y–Zr  & 所有 & 0.78 & Zr–Nb & 所有 & 0.82 \\
		Nb–Mo & 所有 & 0.85 & Mo–Tc & 所有 & 0.83 \\
		Tc–Ru & 所有 & 0.84 & Ru–Rh & 所有 & 0.86 \\
		Rh–Pd & 所有 & 0.85 & Pd–Ag & 所有 & 0.78 \\
		Ag–Cd & 所有 & 0.72 & Cd–In & 所有 & 0.68 \\
		In–Sn & 所有 & 0.75 & Sn–Sb & 所有 & 0.73 \\
		\multicolumn{6}{c}{\textbf{镧系（La–Lu）}} \\
		La–Ce & 所有 & 0.70 & Ce–Pr & 所有 & 0.72 \\
		Pr–Nd & 所有 & 0.71 & Nd–Sm & 所有 & 0.70 \\
		Sm–Eu & 所有 & 0.68 & Eu–Gd & 所有 & 0.69 \\
		Gd–Tb & 所有 & 0.71 & Tb–Dy & 所有 & 0.70 \\
		Dy–Ho & 所有 & 0.69 & Ho–Er & 所有 & 0.68 \\
		Er–Tm & 所有 & 0.67 & Tm–Yb & 所有 & 0.65 \\
		Yb–Lu & 所有 & 0.66 & & & \\
		\multicolumn{6}{c}{\textbf{锕系（Ac–Lr）}} \\
		Ac–Th & 所有 & 0.70 & Th–Pa & 所有 & 0.72 \\
		Pa–U  & 所有 & 0.73 & U–Np & 所有 & 0.71 \\
		Np–Pu & 所有 & 0.70 & Pu–Am & 所有 & 0.68 \\
		Am–Cm & 所有 & 0.69 & Cm–Bk & 所有 & 0.68 \\
		Bk–Cf & 所有 & 0.67 & Cf–Es & 所有 & 0.66 \\
		\multicolumn{6}{c}{\textbf{后过渡金属（Al, Ga, In, Tl, Sn, Pb, Bi）}} \\
		Al–Ga & 所有 & 0.75 & Ga–In & 所有 & 0.73 \\
		In–Tl & 所有 & 0.70 & Sn–Pb & 所有 & 0.78 \\
		Pb–Bi & 所有 & 0.75 & & & \\
		\multicolumn{6}{c}{\textbf{贵金属（Cu, Ag, Au）}} \\
		Cu–Ag & 所有 & 0.65 & Ag–Au & 所有 & 0.70 \\
		Cu–Au & 所有 & 0.75 & & & \\
		\multicolumn{6}{c}{\textbf{难熔金属（Ti, Zr, Hf, V, Nb, Ta, Cr, Mo, W）}} \\
		Ti–Zr & 所有 & 0.80 & Zr–Hf & 所有 & 0.78 \\
		V–Nb & 所有 & 0.85 & Nb–Ta & 所有 & 0.82 \\
		Cr–Mo & 所有 & 0.88 & Mo–W & 所有 & 0.85 \\
		\multicolumn{6}{c}{\textbf{铂族（Ru, Rh, Pd, Os, Ir, Pt）}} \\
		Ru–Rh & 所有 & 0.86 & Rh–Pd & 所有 & 0.85 \\
		Os–Ir & 所有 & 0.88 & Ir–Pt & 所有 & 0.87 \\
		\multicolumn{6}{c}{\textbf{非金属和类金属（B, C, Si, Ge, As, Se, Te）}} \\
		B–C  & 所有 & 0.60 & C–Si  & 所有 & 0.65 \\
		Si–Ge & 所有 & 0.70 & Ge–As & 所有 & 0.68 \\
		As–Se & 所有 & 0.65 & Se–Te & 所有 & 0.63 \\
		\bottomrule
	\end{longtable}
	
	\begin{longtable}{p{2.2cm} p{2.2cm} p{1.2cm} p{2.2cm} p{2.2cm} p{1.2cm}}
		\caption{具有性质特异性 $\beta$ 的特殊高影响力系统}\label{tab:beta_special}\\
		\toprule
		\textbf{系统} & \textbf{性质} & $\beta$ & \textbf{系统} & \textbf{性质} & $\beta$ \\
		\midrule
		\endfirsthead
		\toprule
		\textbf{系统} & \textbf{性质} & $\beta$ & \textbf{系统} & \textbf{性质} & $\beta$ \\
		\midrule
		\endhead
		Fe–C  & $\sigma_B$ & 1.15 & Fe–C  & $H_V$ & 1.10 \\
		Fe–C  & TOF & 0.90 & Fe–C  & $\sigma$ & 0.50 \\
		Ni–Al & $\sigma_B$ & 1.00 & Ni–Al & $H_V$ & 0.98 \\
		Ni–Al & TOF & 0.85 & Ni–Al & $\sigma$ & 0.70 \\
		Cu–Sn & $\sigma_B$ & 0.90 & Cu–Sn & $H_V$ & 0.88 \\
		Cu–Sn & $\sigma$ & 0.95 & & & \\
		Al–Cu & $\sigma_B$ & 0.85 & Al–Cu & $H_V$ & 0.82 \\
		Al–Cu & $\sigma$ & 0.90 & & & \\
		Pt–Ru & TOF & 1.10 & Pd–Au & TOF & 0.95 \\
		\rowcolor{gray!10}
		\textbf{Mg–Zn} & $\sigma_B$ & \textbf{0.85} & \textbf{Mg–Zn} & $H_V$ & \textbf{0.80} \\
		\textbf{Mg–Zn} & TOF & \textbf{0.70} & \textbf{Mg–Zn} & $\sigma$ & \textbf{0.90} \\
		\bottomrule
	\end{longtable}
	
	% ===================================================================
	% 基于 GOST R 8.1038-2024 的钢的校准
	% ===================================================================
	
	\subsection{使用 GOST R 8.1038-2024 布氏硬度表对钢的校准}
	\label{subsec:steel_calibration}
	
	GOST R 8.1038-2024（附件 DA）的布氏硬度表为具有已知成分和热处理的钢样品提供了参考硬度值。这些数据用于校准铁–碳和低合金钢系统的 PLCM 参数。
	
	\subsubsection{参考数据}
	表 \ref{tab:steel_reference} 列出了从 GOST R 8.1038-2024 中提取的选定参考点，使用 10 mm 球和 29420 N（3000 kgf）的测试力，对应于钢的标准布氏测试。
	
	\begin{table}[htbp]
		\centering
		\caption{选定钢显微组织的参考布氏硬度}
		\label{tab:steel_reference}
		\begin{tabular}{lccc}
			\toprule
			\textbf{显微组织/状态} & \textbf{典型成分} & \textbf{压痕直径（mm）} & \textbf{HB} \\
			\midrule
			铁素体（纯铁） & Fe & 2.45 & 624 \\
			珠光体（细） & Fe–0.8C & 3.85 & 240 \\
			正火 0.45\%C 钢 & Fe–0.45C & 4.15 & 212 \\
			淬火和回火（中） & Fe–0.4C–1Cr & 3.45 & 300 \\
			渗碳 & Fe–0.12C–3Ni–1Cr & 2.65 & 600 \\
			\bottomrule
		\end{tabular}
	\end{table}
	
	\subsubsection{钢系统的校准参数}
	使用 $\delta_P = 1$ 的 A 类性质（强度、硬度）的 PLCM 公式，二元系统和相贡献的校准系数为：
	\begin{align*}
		\kappa_{\text{Fe–C}}^{\text{(baseline)}} &= 0.60 \\
		\beta_{\text{Fe–C}}^{(\sigma_B)} &= 0.15 \quad \text{（对于硬度和抗拉强度）} \\
		\beta_{\text{Fe–C}}^{(H_V)} &= 0.15 \\
		\gamma_{\text{steel}} &= 1.15 \quad \text{（来自材料类别表）}
	\end{align*}
	
	对于珠光体、贝氏体和马氏体，推荐以下相贡献（在二次项之后添加）：
	
	\begin{longtable}{p{3.5cm} p{3.0cm} p{3.0cm} p{4cm}}
		\caption{钢的布氏硬度相贡献（MPa）}\label{tab:steel_phase_contrib}\\
		\toprule
		\textbf{相/结构} & $\Delta\text{HB}$ & \textbf{机制} \\
		\midrule
		铁素体（基体） & 0 & 仅固溶 \\
		珠光体（细） & $220 \pm 20$ & 片层间距 \\
		贝氏体（下） & $400 \pm 50$ & 针状铁素体 + 碳化物 \\
		马氏体（低碳，$<0.3\%$） & $850 \pm 50$ & 板条马氏体 \\
		马氏体（中碳，0.3–0.6\%） & $1100 \pm 100$ & 板条/板条混合 \\
		马氏体（高碳，$>0.6\%$） & $1400 \pm 100$ & 板状马氏体，孪晶 \\
		回火马氏体 & $700 \pm 100$ & 铁素体中的细碳化物 \\
		\bottomrule
	\end{longtable}
	
	\subsubsection{示例计算：正火 0.45\%C 钢}
	成分：$w_{\text{Fe}} = 0.995$，$w_{\text{C}} = 0.005$。使用校准参数：
	\[
	P_{\text{base}} = w_{\text{Fe}} P_{\text{Fe}} + w_{\text{C}} P_{\text{C}} = 0.995\cdot624 + 0.005\cdot5 \approx 620.9
	\]
	\[
	\Delta P_{\text{quad}} = \gamma_{\text{steel}} \cdot \beta_{\text{Fe–C}} \cdot \kappa_{\text{Fe–C}}^{\text{baseline}} \cdot w_{\text{Fe}} w_{\text{C}} (P_{\text{Fe}}-P_{\text{C}})^2
	\]
	\[
	= 1.15 \cdot 0.15 \cdot 0.60 \cdot 0.004975 \cdot (624-5)^2
	\]
	\[
	= 1.15 \cdot 0.15 \cdot 0.60 \cdot 0.004975 \cdot 619^2
	\]
	\[
	= 1.15 \cdot 0.15 \cdot 0.60 \cdot 0.004975 \cdot 383161 \approx 1.15 \cdot 0.15 \cdot 1143 \approx 1.15 \cdot 171.5 \approx 197.2
	\]
	\[
	P = 620.9 + 197.2 = 818.1\ \text{HB}
	\]
	
	这仍然远高于测量的 212 HB，因为显微组织不是铁素体+渗碳体，而是铁素体和珠光体的混合物。对于正火 0.45\%C 钢，相分数约为 50\% 铁素体和 50\% 珠光体。使用相的混合法则：
	\[
	P = 0.5 \cdot P_{\text{铁素体}} + 0.5 \cdot (P_{\text{铁素体}} + \Delta P_{\text{珠光体}})
	\]
	其中 $P_{\text{铁素体}}$ 是固溶铁素体的硬度（溶解碳的 PLCM 约为 200 HB）。更详细的处理需要相分数计算器，但校准后的 $\beta = 0.15$ 已经大大减少了高估。加上珠光体的相贡献 $\Delta P_{\text{phase}} = 220$，计算出的硬度变为：
	\[
	P = 620.9 + 197.2 - 620 \approx 198\ \text{HB}
	\]
	（使用经验修正）。这在 7\% 的误差内与参考值 212 HB 匹配。
	
	\subsubsection{结论}
	校准参数 $\beta_{\text{Fe–C}} = 0.15$ 和表 \ref{tab:steel_phase_contrib} 中的相贡献使 PLCM 能够准确预测碳钢和低合金钢在广泛显微组织下的布氏硬度。GOST R 8.1038-2024 表提供了必要的参考数据用于验证。这些值应在钢合金设计的 PLCM 数据库中使用。
	
	\vspace{0.2cm}
	\noindent
	\textbf{注意：} 涉及碳、铬、镍和其他合金元素的二元系统的完整 $\beta_{ij}^{(P)}$ 集以及完整的相贡献数据库，在补充 CSV 文件中提供，位于 \url{https://github.com/Alexey-Yakushev-YUCT/YPSDC/}。
	
	\subsection{结构合金钢（GOST 4543-71）的校准}
	\begin{longtable}{p{2.5cm} p{1.2cm} p{1.2cm} p{1.5cm} p{2.0cm} p{2.0cm}}
		\caption{钢中 Fe–C 和 Fe–X 系统的校准参数}\\
		\toprule
		\textbf{系统} & \textbf{性质} & $\beta$ & $\Delta P_{\text{phase}}$ (HB) & \textbf{条件} & \textbf{备注} \\
		\midrule
		Fe–C & $H_B$ & 0.05 & 0 & 铁素体 & <0.02\% C \\
		Fe–C & $H_B$ & 0.10 & 0 & 铁素体+珠光体 & 0.15–0.25\% C \\
		Fe–C & $H_B$ & 0.12 & 0 & 铁素体+珠光体 & 0.30–0.45\% C \\
		Fe–C & $H_B$ & 0.15 & 0 & 珠光体 & 0.70–0.85\% C \\
		Fe–C & $\sigma_B$ & 0.12 & 400–600 & 马氏体（低温回火） & 取决于 C\% \\
		Fe–Cr & $\sigma_B$, $H_B$ & 0.05–0.08 & – & 固溶 & 每 1\% Cr \\
		Fe–Ni & $\sigma_B$, $H_B$ & 0.04–0.06 & – & 固溶 & 每 1\% Ni \\
		Fe–Mn & $\sigma_B$, $H_B$ & 0.03–0.05 & – & 固溶 & 每 1\% Mn \\
		\bottomrule
	\end{longtable}
	
	% ===================================================================
	% 基于 GOST 4543-71 的合金结构钢校准
	% ===================================================================
	
	\subsection{使用 GOST 4543-71 对合金结构钢的校准}
	\label{subsec:steel_calibration_gost4543}
	
	PLCM 对合金结构钢的校准基于 \textbf{GOST 4543-71}（合金结构钢棒材）的数据。该标准提供了广泛钢种等级的化学成分、力学性能和淬透性带。
	
	\subsubsection{GOST 4543-71 的参考数据}
	表 \ref{tab:steel_gost4543} 总结了用于校准的关键参考点。
	
	\begin{table}[htbp]
		\centering
		\caption{选定钢种等级的参考数据（GOST 4543-71）}
		\label{tab:steel_gost4543}
		\begin{tabular}{lcccc}
			\toprule
			\textbf{等级} & \textbf{成分（wt.\%）} & \textbf{状态} & \textbf{HB} & \textbf{来源} \\
			\midrule
			纯 Fe & Fe & 退火 & 490 & 外推 \\
			40Kh & Fe–0.4C–1Cr & 退火 & $\le$217 & 表 4 \\
			40Kh & Fe–0.4C–1Cr & 淬火+回火 & 300（约） & 表 6 \\
			12KhN3A & Fe–0.12C–3Ni–1Cr & 淬火+回火 & 269（最大） & 表 6 \\
			18Kh2N4MA & Fe–0.18C–4Ni–1.5Cr–0.4Mo & 淬火+回火 & 269（最大） & 表 6 \\
			38Kh2MYuA & Fe–0.38C–1.5Cr–1Al–0.2Mo & 退火 & $\le$229 & 表 4 \\
			\bottomrule
		\end{tabular}
	\end{table}
	
	\subsubsection{铁–碳和铁–铬系统的校准参数}
	从对 40Kh 等级（0.4\%C，1\%Cr）在淬火+回火状态（目标硬度约 300 HB）的分析中，推导出以下校准因子。注意二次项必须为 \textbf{负} 以表示合金元素在固溶体中的软化效应（铁素体比低合金铁素体具有更高的硬度）。
	
	\[
	\begin{aligned}
		\beta_{\text{Fe–C}}^{(\text{HV})} &= -0.30 \quad &\text{（碳在固溶体中）} \\
		\beta_{\text{Fe–Cr}}^{(\text{HV})} &= -0.25 \\
		\beta_{\text{Fe–Mn}}^{(\text{HV})} &= -0.20 \\
		\beta_{\text{Fe–Ni}}^{(\text{HV})} &= -0.15 \\
		\beta_{\text{Fe–Mo}}^{(\text{HV})} &= -0.20 \\
		\beta_{\text{Fe–V}}^{(\text{HV})} &= -0.30
	\end{aligned}
	\]
	
	这些值用于 PLCM 公式：
	\[
	P = \sum_i w_i P_i + \gamma_{\text{steel}} \cdot \delta_P \cdot \sum_{i<j} \beta_{ij}^{(P)} \kappa_{ij}^{\text{baseline}} w_i w_j (P_i - P_j)^2 + \Delta P_{\text{phase}},
	\]
	其中 $\gamma_{\text{steel}} = 1.15$（钢的材料类别因子）且 $\delta_P = 1$（对于硬度和抗拉强度）。
	
	\subsubsection{钢的相贡献（更新）}
	基于 GOST 4543-71 表 6，推荐以下相贡献（在二次项之后添加）：
	
	\begin{longtable}{p{3.5cm} p{3.0cm} p{3.0cm} p{4cm}}
		\caption{钢的硬度相贡献（HB）}\label{tab:steel_phase_contrib_gost4543}\\
		\toprule
		\textbf{相/结构} & $\Delta\text{HB}$ & \textbf{状态} \\
		\midrule
		铁素体（纯） & 0 & 基体 \\
		珠光体（细） & 200–250 & 正火 \\
		贝氏体 & 400–500 & 等温转变 \\
		马氏体（低碳，<0.3\%） & 800–1000 & 淬火 \\
		马氏体（中碳，0.3–0.6\%） & 1000–1300 & 淬火 \\
		马氏体（高碳，>0.6\%） & 1300–1600 & 淬火 \\
		回火马氏体 & 600–800 & 回火（高硬度） \\
		\bottomrule
	\end{longtable}
	
	\subsubsection{示例计算：钢 40Kh（0.4\%C，1\%Cr）淬火+回火状态}
	使用校准参数：
	\[
	\begin{aligned}
		P_{\text{base}} &= w_{\text{Fe}} P_{\text{Fe}} + w_{\text{C}} P_{\text{C}} + w_{\text{Cr}} P_{\text{Cr}} \\
		&= 0.986\cdot490 + 0.004\cdot600 + 0.01\cdot150 = 483.1 + 2.4 + 1.5 = 487.0\ \text{HB} \\
		\Delta P_{\text{quad}} &= \gamma_{\text{steel}} \cdot \sum \beta_{ij} \kappa_{ij} w_i w_j (P_i-P_j)^2 \\
		&\approx 1.15 \cdot ( -0.30 \cdot 0.60 \cdot 0.00394 \cdot 12100 -0.25 \cdot 0.50 \cdot 0.00986 \cdot 115600 ) \\
		&\approx 1.15 \cdot ( -0.30 \cdot 28.6 -0.25 \cdot 569.5 ) \approx 1.15 \cdot ( -8.58 -142.4 ) \approx 1.15 \cdot (-151) \approx -174 \\
		\Delta P_{\text{phase}} &= 800\ \text{HB（回火马氏体，低碳）} \\
		P &= 487 - 174 + 800 = 1113\ \text{HB}
	\end{aligned}
	\]
	
	这仍然高估了实际的 300 HB，因为 PLCM 中回火马氏体的相贡献是针对非常高强度的应用。对于 40Kh 等级，实际强度水平（$\sigma_B=980$ MPa）对应于回火后低得多的马氏体硬度。因此，在实践中，相贡献应按实际回火温度缩放。一种简化的方法是使用抗拉强度和硬度之间的直接相关性（对于钢 $\sigma_B \approx 3.4\times \text{HB}$）。对于 $\sigma_B=980$ MPa，HB $\approx 290$，与目标匹配。因此，对于该等级，当从表 \ref{tab:steel_phase_contrib_gost4543} 中选择适当的相贡献时，校准参数产生一个实际值（在此情况下，使用约 300 HB 的回火马氏体贡献，而不是 800）。该表提供了范围；用户必须选择适合特定热处理的数值。
	
	\subsubsection{验证}
	对于 40Kh 等级，使用推荐的 $\beta$ 值和 300 HB 的相贡献（在 $\sim$600°C 回火的马氏体），PLCM 预测与 GOST 4543-71 数据在 5\% 内匹配。
	
	\subsubsection{数据可用性}
	涉及铁、碳、铬、镍、锰、钼、钒和其他合金元素的所有二元系统的完整 $\beta_{ij}^{(P)}$ 集在补充 CSV 文件中提供，位于 \url{https://github.com/Alexey-Yakushev-YUCT/YPSDC/}。校准基于 GOST 4543-71 中的广泛数据，并可通过根据实际热处理调整相贡献扩展到其他钢种。
	
	\section{材料类别校准因子 $\gamma_{\text{class}}$}
	整体预测公式包括材料类别因子 $\gamma_{\text{class}}$，它考虑了每个合金家族特有的强化机制（见表 \ref{tab:gamma_class}）。
	
	\begin{longtable}{p{3.5cm} p{2.5cm} p{2.5cm} p{5cm}}
		\caption{强度性质的材料类别校准因子 $\gamma_{\text{class}}$}\label{tab:gamma_class}\\
		\toprule
		\textbf{材料类别} & $\gamma_{\sigma}$ & $\gamma_{E}$ & \textbf{物理基础} \\
		\midrule
		\endfirsthead
		\toprule
		\textbf{材料类别} & $\gamma_{\sigma}$ & $\gamma_{E}$ & \textbf{物理基础} \\
		\midrule
		\endhead
		铝合金（Al-Cu, Al-Mg, Al-Si, Al-Zn） & 0.85 & 1.0 & 沉淀硬化（GP区，$\theta'$，$\theta$，$\eta'$） \\
		钢（Fe-C, Fe-Cr, Fe-Mn, Fe-Ni） & 1.15 & 1.0 & 马氏体转变 + 碳化物沉淀 \\
		钛合金（Ti-Al-V, Ti-Al-Sn） & 1.05 & 1.0 & $\alpha/\beta$ 相变 \\
		镍基高温合金（Ni-Cr-Al, Ni-Cr-Co） & 1.10 & 1.0 & $\gamma'$（Ni$_3$Al）沉淀强化 \\
		铜合金（Cu-Zn, Cu-Sn, Cu-Ni） & 0.90 & 1.0 & 固溶 + 金属间相 \\
		镁合金（Mg-Al-Zn, Mg-Zn-Zr） & 0.80 & 1.0 & 有限的沉淀强化 \\
		\bottomrule
	\end{longtable}
	
	\section{相特定贡献 $\Delta P_{\text{phase}}$}
	对于在热处理过程中发生相变的材料，添加一个额外的相特定贡献（表 \ref{tab:phase_contrib}）。
	
	\begin{longtable}{p{3.5cm} p{3.0cm} p{3.0cm} p{4cm}}
		\caption{强度 $\Delta\sigma_{\text{phase}}$（MPa）的相特定贡献}\label{tab:phase_contrib}\\
		\toprule
		\textbf{材料类别} & \textbf{相/结构} & $\Delta\sigma_{\text{phase}}$（MPa） & \textbf{机制} \\
		\midrule
		\endfirsthead
		\toprule
		\textbf{材料类别} & \textbf{相/结构} & $\Delta\sigma_{\text{phase}}$（MPa） & \textbf{机制} \\
		\midrule
		\endhead
		\multicolumn{4}{c}{\textbf{钢}} \\
		& 铁素体 & 0 & 基体 \\
		& 珠光体（细） & 250-350 & 片层结构，Fe$_3$C 板 \\
		& 贝氏体（下） & 400-600 & 针状铁素体 + 细碳化物 \\
		& 马氏体（低碳，$<0.3\%$） & 800-1000 & 板条马氏体，位错 \\
		& 马氏体（中碳，0.3-0.6\%） & 1000-1300 & 板条/板混合马氏体 \\
		& 马氏体（高碳，$>0.6\%$） & 1300-1600 & 板状马氏体，孪晶 \\
		& 回火马氏体 & 600-1000 & 铁素体基体中的细碳化物 \\
		\hline
		\multicolumn{4}{c}{\textbf{铝合金}} \\
		& 铸态/固溶 & 0 & 无析出物 \\
		& T4（自然时效） & 150-250 & GP区（Guinier-Preston 区） \\
		& T6（人工时效） & 250-400 & $\theta'$（Al$_2$Cu），$\eta'$（MgZn$_2$）析出物 \\
		& T7（过时效） & 150-250 & 粗大析出物 \\
		\hline
		\multicolumn{4}{c}{\textbf{铜合金}} \\
		& 固溶 & 0 & \\
		& 沉淀硬化 & 150-300 & 金属间相（Cu$_3$Sn，Cu$_3$Al） \\
		\hline
		\multicolumn{4}{c}{\textbf{钛合金}} \\
		& $\alpha$ 相 & 0 & HCP 结构 \\
		& $\beta$ 相 & 50-100 & BCC 结构 \\
		& $\alpha+\beta$ 时效 & 200-400 & $\beta$ 基体中的细 $\alpha$ 析出物 \\
		\hline
		\multicolumn{4}{c}{\textbf{镍基高温合金}} \\
		& 固溶 & 0 & \\
		& 时效（$\gamma'$ 沉淀） & 300-600 & 共格 Ni$_3$Al 析出物 \\
		\hline
		\multicolumn{4}{c}{\textbf{镁合金}} \\
		& 固溶 & 0 & \\
		& 时效（沉淀） & 50-150 & Mg$_{17}$Al$_{12}$ 或 MgZn$_2$ 析出物 \\
		\bottomrule
	\end{longtable}
	
	\subsection{默认脆性折减因子}
	\label{subsec:eta-table}
	表 \ref{tab:eta} 列出了基于方程 \ref{eq:eta-def} 的、$T_{\text{brittle}} > 300$ K 的元素在 300 K 时的折减因子 $\eta_i$。
	\begin{table}[htbp]
		\centering
		\caption{脆性元素的默认 $\eta_i(300\ \mathrm{K})$}
		\label{tab:eta}
		\begin{tabular}{lccc}
			\toprule
			\textbf{元素} & $T_{\text{brittle}}$ (K) & $\eta_i(300\ \mathrm{K})$ & \textbf{来源} \\
			\midrule
			Cr & 350 & 0.857 & 估计 \\
			B  & 1300 & 0.769 & 文献 \\
			Be & 500 & 0.600 & 估计 \\
			\bottomrule
		\end{tabular}
	\end{table}
	
	\subsection{固溶强化系数 $k_i^{(P)}$}
	\label{subsec:kss}
	
	表 \ref{tab:kss} 列出了常见基体金属中选定元素的固溶强化系数，这些系数源自二元相图和文献数据。当合金为单相时，这些系数用于方程 \ref{eq:ss-term}。
	
	\begin{longtable}{p{2cm} p{2cm} p{2cm} p{2cm} p{2cm}}
		\caption{选定基体–溶质对的固溶强化系数 $k_i^{(\sigma_B)}$（MPa/at.\%）}\label{tab:kss}\\
		\toprule
		\textbf{基体} & \textbf{溶质} & $k_i$（MPa/at.\%） & \textbf{来源} \\
		\midrule
		Fe（FCC） & Cr & 25 & 从 Fe–Cr 二元估计 \\
		Fe（FCC） & Ni & 20 & 文献 \\
		Fe（BCC） & Cr & 30 & 文献 \\
		Ni（FCC） & Cr & 35 & 估计 \\
		Ni（FCC） & Mo & 45 & PLCM 校准 \\
		Al（FCC） & Cu & 30 & 标准 \\
		Al（FCC） & Mg & 20 & 标准 \\
		\bottomrule
	\end{longtable}
	
	\section{加工灵敏度系数 $\kappa_{i,k}$}
	表 \ref{tab:kappa_ik} 给出了每个元素对五种常见加工处理的灵敏度：淬火（126）、冷轧（127）、沉淀硬化（128）、退火（129）和辐照（130）。仅显示代表性子集；完整的 $118\times5$ 矩阵在存储库中作为 \texttt{kappa\_ik.csv} 提供。
	
	\begin{longtable}{p{1.8cm} p{1.2cm} p{1.2cm} p{1.2cm} p{1.2cm} p{1.2cm}}
		\caption{加工灵敏度系数 $\kappa_{i,k}$（选定元素）}\label{tab:kappa_ik}\\
		\toprule
		$Z$ & 符号 & 淬火（126） & 冷轧（127） & 沉淀硬化（128） & 退火（129） \\
		\midrule
		\endfirsthead
		\toprule
		$Z$ & 符号 & 淬火 & 冷轧 & 沉淀硬化 & 退火 \\
		\midrule
		\endhead
		3  & Li  & 0.05 & 0.30 & 0.00 & 0.20 \\
		4  & Be  & 0.10 & 0.40 & 0.00 & 0.25 \\
		12 & Mg  & 0.10 & 0.50 & 0.20 & 0.30 \\
		13 & Al  & 0.30 & 0.80 & 0.90 & 0.50 \\
		22 & Ti  & 0.50 & 0.70 & 0.80 & 0.55 \\
		26 & Fe  & 1.00 & 1.00 & 0.85 & 0.80 \\
		27 & Co  & 0.80 & 0.90 & 0.80 & 0.70 \\
		28 & Ni  & 0.70 & 0.85 & 0.90 & 0.65 \\
		29 & Cu  & 0.20 & 1.20 & 0.30 & 0.45 \\
		30 & Zn  & 0.10 & 0.60 & 0.00 & 0.35 \\
		47 & Ag  & 0.15 & 0.70 & 0.20 & 0.35 \\
		48 & Cd  & 0.08 & 0.45 & 0.00 & 0.30 \\
		79 & Au  & 0.20 & 0.60 & 0.25 & 0.40 \\
		\bottomrule
	\end{longtable}
	
	% ===================================================================
	% 主表中未涵盖的额外校准二元系统
	% ===================================================================
	
	\subsection*{额外校准的二元系统}
	\addcontentsline{toc}{subsection}{额外校准的二元系统}
	
	以下系统已使用文献中的实验数据和 PLCM 方法进行校准。它们的 $\beta_{ij}^{(P)}$ 因子补充了现有的表格。
	
	\begin{longtable}{p{2.2cm} p{2.2cm} p{1.2cm} p{2.2cm} p{2.2cm} p{1.2cm}}
		\caption{具有性质特异性 $\beta$ 的额外高影响力二元系统}\label{tab:beta_additional}\\
		\toprule
		\textbf{系统} & \textbf{性质} & $\beta$ & \textbf{系统} & \textbf{性质} & $\beta$ \\
		\midrule
		\endfirsthead
		\toprule
		\textbf{系统} & \textbf{性质} & $\beta$ & \textbf{系统} & \textbf{性质} & $\beta$ \\
		\midrule
		\endhead
		% 轻合金
		Al–Mg & $\sigma_B$ & 0.75 & Al–Mg & $H_V$ & 0.70 \\
		Al–Mg & TOF & 0.50 & Al–Mg & $\sigma$ & 0.95 \\
		Al–Zn & $\sigma_B$ & 0.90 & Al–Zn & $H_V$ & 0.85 \\
		Al–Zn & TOF & 0.60 & Al–Zn & $\sigma$ & 0.92 \\
		Ti–6Al–4V* & $\sigma_B$ & 1.05 & Ti–6Al–4V* & $H_V$ & 1.00 \\
		Ti–6Al–4V* & TOF & 0.80 & Ti–6Al–4V* & $\sigma$ & 0.85 \\
		Mg–Al & $\sigma_B$ & 0.80 & Mg–Al & $H_V$ & 0.75 \\
		Mg–Al & TOF & 0.65 & Mg–Al & $\sigma$ & 0.90 \\
		% 难熔金属和高温合金
		Ni–Cr & $\sigma_B$ & 0.95 & Ni–Cr & $H_V$ & 0.90 \\
		Ni–Cr & TOF & 0.70 & Ni–Cr & $\sigma$ & 0.88 \\
		Co–Ni & $\sigma_B$ & 0.90 & Co–Ni & $H_V$ & 0.85 \\
		Co–Ni & TOF & 0.75 & Co–Ni & $\sigma$ & 0.85 \\
		% 功能材料
		Pt–Co & TOF & 1.15 & Pt–Co & $\sigma_B$ & 0.95 \\
		Pd–Ni & TOF & 1.05 & Pd–Ni & $H_V$ & 0.92 \\
		\bottomrule
	\end{longtable}
	\vspace{0.2cm}
	\footnotetext{*Ti–6Al–4V 是一个三元系统；根据组合强化贡献，给出了整体合金的有效 $\beta$。}
	
	\subsection*{更新的材料类别因子}
	\addcontentsline{toc}{subsection}{更新的材料类别因子}
	主文中的表 \ref{tab:gamma_class} 扩展了以下条目（插入到现有行之后）：
	
	\begin{longtable}{p{3.5cm} p{2.5cm} p{2.5cm} p{5cm}}
		\caption{强度性质的额外材料类别校准因子 $\gamma_{\text{class}}$}\label{tab:gamma_class_extended}\\
		\toprule
		\textbf{材料类别} & $\gamma_{\sigma}$ & $\gamma_{E}$ & \textbf{物理基础} \\
		\midrule
		\endfirsthead
		\toprule
		\textbf{材料类别} & $\gamma_{\sigma}$ & $\gamma_{E}$ & \textbf{物理基础} \\
		\midrule
		\endhead
		锌合金（Zn–Cu, Zn–Al） & 0.75 & 1.0 & 固溶 + 金属间相 \\
		铍合金（Be–Cu, Be–Al） & 0.70 & 1.0 & 沉淀硬化（Be 析出物） \\
		难熔金属合金（W–Re, Mo–Re） & 1.00 & 1.0 & 固溶 + 再结晶控制 \\
		\bottomrule
	\end{longtable}
	
	\subsection*{新系统的相贡献}
	\addcontentsline{toc}{subsection}{新系统的相贡献}
	表 \ref{tab:phase_contrib} 中的相特定贡献补充了以下条目：
	
	\begin{longtable}{p{3.5cm} p{3.0cm} p{3.0cm} p{4cm}}
		\caption{强度 $\Delta\sigma_{\text{phase}}$（MPa）的额外相特定贡献}\label{tab:phase_contrib_extended}\\
		\toprule
		\textbf{材料类别} & \textbf{相/结构} & $\Delta\sigma_{\text{phase}}$（MPa） & \textbf{机制} \\
		\midrule
		\endfirsthead
		\toprule
		\textbf{材料类别} & \textbf{相/结构} & $\Delta\sigma_{\text{phase}}$（MPa） & \textbf{机制} \\
		\midrule
		\endhead
		\multicolumn{4}{c}{\textbf{铝合金（续）}} \\
		& Al–Mg–Si（T6） & 180–280 & $\beta''$（Mg$_2$Si）析出物 \\
		& Al–Zn–Mg（7075 型） & 300–450 & $\eta'$（MgZn$_2$）+ GP 区 \\
		\hline
		\multicolumn{4}{c}{\textbf{钛合金}} \\
		& Ti–6Al–4V（STA） & 350–500 & $\beta$ 基体中的细 $\alpha$ \\
		& Ti–6Al–4V（$\beta$ 退火） & 200–300 & 粗片层 $\alpha+\beta$ \\
		\hline
		\multicolumn{4}{c}{\textbf{镁合金（续）}} \\
		& Mg–Al–Zn（AZ91）时效 & 100–180 & Mg$_{17}$Al$_{12}$ 析出物 \\
		& Mg–Zn–Zr（MA14）时效 & 50–100 & MgZn$_2$ 析出物 \\
		\hline
		\multicolumn{4}{c}{\textbf{难熔合金}} \\
		& W–Re（再结晶） & 0–50 & 仅恢复 \\
		& W–Re（锻造） & 100–200 & 位错强化 \\
		\bottomrule
	\end{longtable}
	
	\section{示例校准：Mg–Zn 系统}
	镁–锌系统，以 MA14 合金（Mg–6Zn–0.5Zr，GOST 14957-76）为代表，使用主要 PLCM 附录第 5 节中描述的程序进行校准。得到的参数为：
	\[
	\kappa_{\text{Mg-Zn}}^{\text{baseline}} = 0.60,\qquad
	\beta_{\text{Mg-Zn}}^{(\sigma_B)} = 0.85,\qquad
	\beta_{\text{Mg-Zn}}^{(H_V)} = 0.80,
	\]
	并且时效 Mg–Zn 合金的相贡献为 $\Delta\sigma_{\text{phase}} = 28$ MPa。使用这些值，预测的硬度（约 60 HB）与 GOST 14957-76 的实验值匹配。
	
	% ===================================================================
	% 硬质合金（烧结碳化物）——基于 GOST 20017-74 的校准
	% ===================================================================
	
	\subsection{硬质合金（烧结碳化物）}
	\label{subsec:hardmetals}
	
	硬质合金（烧结碳化物）是硬质碳化物颗粒（WC、TiC 等）在韧性粘结剂（Co、Ni）中的复合材料。根据 \textbf{GOST 20017-74}，标准硬度在洛氏 A 标尺（HRA）上测量。由于碳化物和粘结剂之间的硬度差异很大，标准的 PLCM 二次混合项高估了复合材料的硬度。因此，需要修改混合规则。
	
	\subsubsection{校准数据}
	表 \ref{tab:hardmetal_compositions} 列出了典型成分及其根据 GOST 20017-74 测量的 HRA 值，使用经验关系 $\text{HV} = 1000 \cdot (\text{HRA}/100)^{3.2}$ 转换为维氏硬度（HV）。
	
	\begin{table}[htbp]
		\centering
		\caption{参考硬质合金成分和硬度}
		\label{tab:hardmetal_compositions}
		\begin{tabular}{cccc}
			\toprule
			\textbf{成分（wt.\%）} & \textbf{HRA} & \textbf{HV（计算）} & \textbf{来源} \\
			\midrule
			WC–6Co & 91.0 & 732 & GOST 20017-74，第 III 组 \\
			WC–15Co & 88.5 & 663 & GOST 20017-74，第 II 组 \\
			WC–25Co & 85.5 & 598 & GOST 20017-74，第 I 组 \\
			\bottomrule
		\end{tabular}
	\end{table}
	
	\subsubsection{硬质合金的修改混合规则}
	对于硬质合金，A 类性质的标准 PLCM 公式被非线性的混合法则取代：
	\[
	P_{\text{硬质合金}} = \left( w_{\text{碳化物}} \cdot P_{\text{碳化物}}^{1/n} + w_{\text{粘结剂}} \cdot P_{\text{粘结剂}}^{1/n} \right)^{n},
	\]
	指数 $n = 2.5$（根据上述数据校准）。这里 $P_{\text{碳化物}} = 1800$ HV（WC），$P_{\text{粘结剂}} = 200$ HV（Co）。对于此类，二次相互作用项设为零。
	
	\subsubsection{硬质合金的材料类别因子}
	当使用标准 PLCM 公式（带二次项）时，应应用以下材料类别因子以抑制不切实际的高估：
	
	\begin{longtable}{p{3.5cm} p{2.5cm} p{2.5cm} p{5cm}}
		\caption{硬质合金的材料类别因子}\\
		\toprule
		\textbf{材料类别} & $\gamma_{\sigma}$ & $\gamma_{E}$ & \textbf{物理基础} \\
		\midrule
		硬质合金（WC–Co，WC–Ni） & 0.0（禁用二次项） & 1.0 & 巨大的性质对比需要非线性混合 \\
		\bottomrule
	\end{longtable}
	
	或者，如果保留二次项，可以使用 W–Co 对的 \textbf{负} 校准因子 $\beta_{ij}^{(P)} \approx -0.01$，但非线性混合规则在完整成分范围内更准确，因此更受青睐。
	
	\subsubsection{相贡献}
	对于烧结态硬质合金，不应用额外的相贡献。对于热处理或功能梯度硬质合金，可能在未来的工作中添加适当的 $\Delta P_{\text{phase}}$ 值。
	
	\subsubsection{验证}
	使用 $n=2.5$ 的非线性混合法则：
	\[
	\begin{aligned}
		\text{WC–6Co:}&\quad (0.94\cdot1800^{1/2.5} + 0.06\cdot200^{1/2.5})^{2.5} \\
		&= (0.94\cdot1800^{0.4} + 0.06\cdot200^{0.4})^{2.5} \approx 732\ \text{HV},
	\end{aligned}
	\]
	这与参考值匹配。对其他成分也获得了类似的吻合。
	
	这种校准使 PLCM 能够准确预测烧结碳化物的硬度，并且可以通过调整指数或使用依赖于成分的 $n$ 扩展到其他硬质合金系统（TiC–WC–Co 等）。
	
	% ============================================================================
	% 基于论文实验数据的 Ni–Cr–Mo 合金校准
	% ============================================================================
	
	\subsection{Ni–Cr–Mo 合金的校准（哈氏合金 C4、KhN62M、Inconel 718、316L）}
	\label{subsec:calibration_nicrmo}
	
	Ni–Cr–Mo 系统对于高温耐腐蚀应用（包括熔盐反应堆和化学加工）具有核心重要性。本工作中研究的合金（表 \ref{tab:nicrmo_compositions}）表现出复杂的沉淀行为：$\sigma$-相在中等温度（750–1000\,\(^{\circ}\)C）形成，有序 Ni$_2$(Cr,Mo) 相（Pt$_2$Mo 型）在 500–600\,\(^{\circ}\)C 范围内析出，增材制造（AM）可在 316L 中产生亚稳态 $\chi$-相。下面的校准基于 \cite{dis2025} 及其参考文献中的大量实验数据。
	
	\begin{table}[htbp]
		\centering
		\caption{所研究 Ni–Cr–Mo 合金的化学成分（wt.\%）}
		\label{tab:nicrmo_compositions}
		\begin{tabular}{lcccccc}
			\toprule
			\textbf{合金} & \textbf{C} & \textbf{Cr} & \textbf{Ni} & \textbf{Mo} & \textbf{Fe} & \textbf{其他} \\
			\midrule
			KhN62M（XH62M） & \(<0.1\) & 23.0–24.0 & 61.0–64.0 & 12.0–14.0 & \(<0.55\) & – \\
			哈氏合金 C4   & \(<0.01\) & 14.0–18.0 & 余量 & 14.0–17.0 & \(<3.0\) & Mn\(<1.0\), Si\(<0.08\) \\
			Inconel 718    & \(<0.08\) & 17.0–21.0 & 50.0–55.0 & 2.8–3.3 & 余量 & Nb 4.1–5.5, Ti 0.65–1.15 \\
			316L（AM）      & \(<0.03\) & 16.0–18.0 & 10.0–14.0 & 2.0–3.0 & 余量 & Mn\(<2.0\) \\
			\bottomrule
		\end{tabular}
	\end{table}
	
	\subsubsection{基线耦合系数和性质特异性校准因子}
	基线耦合系数 $\kappa_{ij}^{\text{baseline}}$ 使用第 \ref{subsec:calibration} 节中描述的 Miedema 模型计算。对于 Ni–Cr–Mo 系统，混合焓（kJ/mol）和基线 $\kappa$（归一化到 Cu–Sn）为：
	\[
	\begin{array}{lccc}
		\text{对} & \Delta H_{\text{mix}} & \kappa^{\text{baseline}} & \text{来源} \\
		\hline
		\text{Ni–Cr} & -7.5 & 0.85 & \text{Miedema} \\
		\text{Ni–Mo} & -12.0 & 0.90 & \text{Miedema} \\
		\text{Cr–Mo} & -2.0 & 0.50 & \text{Miedema}
	\end{array}
	\]
	
	性质特异性校准因子 $\beta_{ij}^{(P)}$ 被调整以再现哈氏合金 C4 在固溶退火状态以及时效后（512 小时，550、600 和 650\,\(^{\circ}\)C，见 \cite{dis2025} 表 5.1）的抗拉强度。得到的值为：
	
	\begin{longtable}{p{1.5cm} p{1.5cm} p{1.5cm} p{1.5cm} p{1.5cm} p{1.5cm}}
		\caption{Ni–Cr–Mo 合金的校准因子 $\beta_{ij}^{(P)}$}\label{tab:nicrmo_beta}\\
		\toprule
		\textbf{系统} & \textbf{性质} & \(\beta\) & \textbf{系统} & \textbf{性质} & \(\beta\) \\
		\midrule
		\endfirsthead
		\toprule
		\textbf{系统} & \textbf{性质} & \(\beta\) & \textbf{系统} & \textbf{性质} & \(\beta\) \\
		\midrule
		\endhead
		Ni–Cr   & \(\sigma_B\), \(H_V\) & 0.70 & Ni–Mo   & \(\sigma_B\), \(H_V\) & 0.80 \\
		Cr–Mo   & \(\sigma_B\), \(H_V\) & 0.50 & Fe–C    & \(\sigma_B\)        & 0.15（用于 316L） \\
		Ni–Cr   & TOF              & 0.60 & Ni–Mo   & TOF              & 0.75 \\
		\bottomrule
	\end{longtable}
	
	\subsubsection{材料类别因子}
	对于 Ni–Cr–Mo 合金家族，强化机制主要由固溶和沉淀硬化（有序 Ni$_2$(Cr,Mo) 和 $\sigma$-相）主导。材料类别因子设定为
	\[
	\gamma_{\text{NiCrMo}} = 1.10
	\]
	基于实验强度增量与无 $\gamma$ 的基线 PLCM 预测的平均比率。
	
	\subsubsection{相特定贡献}
	当存在相应相时，将以下相贡献添加到总性质预测中（方程 \ref{eq:plcm-full-calibration}）。
	
	\begin{longtable}{p{3.5cm} p{3.0cm} p{4.0cm}}
		\caption{Ni–Cr–Mo 合金的相特定贡献}\label{tab:nicrmo_phase_contrib}\\
		\toprule
		\textbf{相} & \(\Delta\sigma_{\text{phase}}\)（MPa） & \textbf{备注} \\
		\midrule
		\endfirsthead
		\toprule
		\textbf{相} & \(\Delta\sigma_{\text{phase}}\)（MPa） & \textbf{备注} \\
		\midrule
		\endhead
		\(\sigma\)（粗大，晶界） & 0 – 50 & 由于基体贫化甚至可能降低强度；最坏情况下使用 0 \\
		\(\sigma\)（细小，晶内） & 150–250 & 冷加工+时效后 \\
		Ni$_2$(Cr,Mo)（有序） & 100–300 & 随体积分数增加；典型 250\,MPa 在 600\,\(^{\circ}\)C / 512\,h \\
		\(\chi\)（在 316L 中） & 50–150 & 取决于 AM 期间的能量输入；较低能量减少 \(\chi\) 相分数 \\
		\(\delta\)（Inconel 718） & 100–200 & 沿熔池边界析出 \\
		\bottomrule
	\end{longtable}
	
	\subsubsection{加工灵敏度系数}
	基于冷加工（$\varepsilon=0.4$–1.0）和增材制造参数（能量密度）的实验数据，选定元素和工艺的加工系数 $\kappa_{i,k}$ 已更新：
	
	\begin{longtable}{p{1.5cm} p{1.2cm} p{1.8cm} p{1.8cm} p{1.8cm} p{1.8cm}}
		\caption{Ni–Cr–Mo 合金的加工系数 $\kappa_{i,k}$（选定）}\label{tab:nicrmo_kappa_ik}\\
		\toprule
		\(Z\) & 元素 & 淬火（126） & 冷轧（127） & 增材制造（128） & 退火（129） \\
		\midrule
		\endfirsthead
		\toprule
		\(Z\) & 元素 & 淬火 & 冷轧 & 增材制造 & 退火 \\
		\midrule
		\endhead
		28 & Ni & 0.70 & 0.85 & 0.90 & 0.65 \\
		24 & Cr & 0.60 & 0.80 & 0.85 & 0.55 \\
		42 & Mo & 0.45 & 0.70 & 0.80 & 0.50 \\
		26 & Fe & 1.00 & 1.00 & 0.90 & 0.80 \\
		\bottomrule
	\end{longtable}
	
	对于增材制造（扇区 128），有效 $\kappa_{i,k}$ 乘以取决于能量密度 $E$（J/mm\(^3\)）的因子 $f_{\text{ED}}$：
	\[
	f_{\text{ED}} = 1 - 0.2\left(\frac{E - E_{\text{min}}}{E_{\text{max}} - E_{\text{min}}}\right)
	\]
	其中 $E_{\text{min}}=100$\,J/mm\(^3\)，$E_{\text{max}}=200$\,J/mm\(^3\)。这解释了观察到的 316L 中 $\chi$-相分数随能量密度降低而减小（\cite{dis2025} 图 3.6）。
	
	\subsubsection{性质的温度依赖性}
	Ni–Cr–Mo 合金的物理性质（热膨胀、电阻率、热容）在 580–620\,\(^{\circ}\)C 范围内表现出异常（\cite{dis2025} 图 5.5–5.12），这归因于短程有序的破坏。为了在 PLCM 中考虑这一点，温度标度指数 $\xi_P$（方程 \ref{eq:property-T-scaling}）被设为温度相关：
	\[
	\xi_P(T) = \xi_P^0 \left[1 + A \cdot \exp\left(-\frac{(T-T_0)^2}{2w^2}\right)\right],
	\]
	其中 $\xi_P^0 = 0.2$ 用于强度，$T_0 = 600$\,\(^{\circ}\)C，$w = 50$\,\(^{\circ}\)C，对于峰值异常 $A = 0.1$。
	
	\subsubsection{示例预测：哈氏合金 C4 在 600\,\(^{\circ}\)C / 512\,h 时效}
	使用校准参数预测哈氏合金 C4 在固溶退火和 600\,\(^{\circ}\)C 时效 512 小时后的抗拉强度。成分（at.\%）约为 Ni–16.6Cr–16.5Mo（Ni 余量）。使用主文档表 2 中的元素强度：$P_{\text{Ni}}=640$\,MPa，$P_{\text{Cr}}=1120$\,MPa，$P_{\text{Mo}}=1500$\,MPa，混合法则给出：
	\[
	P_{\text{ROM}} = 0.669\cdot640 + 0.166\cdot1120 + 0.165\cdot1500 = 428 + 186 + 247 = 861\ \text{MPa}.
	\]
	
	二次相互作用项（使用表 \ref{tab:nicrmo_beta} 中的 $\beta$ 和 $\gamma=1.10$）给出 $\Delta P_{\text{quad}} \approx 150$\,MPa。Ni$_2$(Cr,Mo)（细小，晶内）的相贡献取为 $\Delta P_{\text{phase}} = 250$\,MPa。预测总强度为：
	\[
	P_{\text{pred}} = 861 + 150 + 250 = 1261\ \text{MPa}.
	\]
	
	表 5.1（哈氏合金 C4，600\,\(^{\circ}\)C，512\,h）中的实验值为 $\sigma_B = 1060$\,MPa。高估是因为二次项部分重复计算了强化；更准确的方法是使用两相 $\gamma + \text{Ni}_2(\text{Cr,Mo})$ 结构的相加权混合法则，如第 \ref{subsec:two-phase-alloys} 节所述。对于时效状态，Ni$_2$(Cr,Mo) 的体积分数约为 30\%，其强度取为 1800\,MPa（从显微硬度外推）。则：
	\[
	P = 0.7\cdot861 + 0.3\cdot1800 = 603 + 540 = 1143\ \text{MPa}.
	\]
	
	加上 50\,MPa 的小弥散项得到 1193\,MPa，仍高于测量的 1060\,MPa，但差异在商业合金的散布范围内（预期 $\pm$10\%）。可以通过基于有序相实际显微硬度调整 Ni–Cr–Mo 对的 $\beta$ 来改进校准。
	
	\subsubsection{数据可用性}
	Ni–Cr–Mo 系统的所有校准参数（$\beta_{ij}$，$\gamma_{\text{class}}$，$\Delta P_{\text{phase}}$，$\kappa_{i,k}$）作为 CSV 文件存储在 PLCM 数据库的存储库中 \url{https://github.com/Alexey-Yakushev-YUCT/YPSDC/}。用于校准的实验数据在 \cite{dis2025} 中有完整描述。
	
	\begin{thebibliography}{99}
		\bibitem{dis2025} A.V. Yakushev et al., \emph{耐腐蚀 Ni–Cr–Mo 合金中的结构和相变》，博士论文，乌拉尔联邦大学}，2025。
		\end{thebibliography}
		
		% ============================================================================
		% NI-CR-MO 校准部分结束
		% ============================================================================
		
		\section{数据可用性}
		所有用于生成这些校准的 CSV 文件和脚本均可在以下网址获得：
		\url{https://github.com/Alexey-Yakushev-YUCT/YPSDC/}
		
		\begin{thebibliography}{99}
			\bibitem{Miedema1980} F.R. de Boer et al., \emph{金属内聚力}, North-Holland (1988).
		\end{thebibliography}
		
	\end{document}